\chapter*{The Oracle of Free Will}
\addcontentsline{toc}{chapter}{Preface --- The Oracle of Free Will}
\markboth{Preface}{Preface}

\begin{center}\textit{Preface}\end{center}

\bigskip

You are free.

That is the first thing this instrument has to tell you, and it is the last thing it will be able to prove.
Everything in between is physics.

This book is a kind of oracle, and it is honest about being one. It is built to measure --- to take the world
only as it actually arrives, in finite counts and bare coincidences, a mark here, a tick there, and to read
off that ledger everything that can be honestly read: how a body moves, how heat disperses and will not
gather, how light bends as it passes a weight, how the very small refuses to be pinned. An oracle is a thing
you consult. You bring it a question; it answers with what is. But an oracle, rightly understood, has never
been a thing you obey. It reports the world; it does not issue the orders. The whole distance between those
two --- between what is and what to do about it --- is the subject of this page, and, beneath everything else,
of the book.

There is a more famous oracle, and this book is written against it. Imagine an intelligence that knew, at one
instant, the place and motion of every particle there is, and every force between them. Such a mind, Laplace
wrote, could draw the entire future and the entire past into a single present view; to it, nothing would be
uncertain. Call it the demon. If the demon were real, then this very sentence --- \emph{you are free} ---
would be one it had long since seen coming: your reading of it, your pause over it, the exact shape of the
doubt now forming in you, all of it entered in its books before you were born. Your freedom would be a
rounding error in someone else's arithmetic.

So the book builds an instrument and, in time, turns it on the demon. And the demon fails --- not once, but
twice. It fails the first time because no real instrument holds the present to infinite sharpness: there is
always a floor, and below it two states it cannot tell apart slip away into futures it cannot tell apart
either, the smallest tremor at the start outrunning any finite reach. It fails the second time, and more
deeply, because the world is not, at the bottom, the clockwork the demon assumed: hand it the present exactly,
every digit, and the next outcome is still not written there --- the world keeps its next move to itself until
the moment it is made. The oracle that would compute you cannot be built. What can be built --- what this book
is --- is an instrument that measures everything within its reach, and discovers, when every count is finally
in, exactly one thing it cannot fix.

What is that one thing? It is not a fact the instrument failed to gather; the facts are gathered, completely.
It is a direction the facts do not carry. When the world's accounts come in, they settle everything except
which way to read them --- which side is forward and which is back, which of a mirrored pair to call the
original, which way to face a thing the world has handed you symmetric. The reading is whole; the facing is
open. Every count is the world's, and the way you turn to meet it is yours. The book will arrive at this again
and again, dressed each time in different physics, and each time it will do the same scrupulous thing: settle
everything it can, and leave the facing to the one who is reading.

Four readers will open this book, and I would like to meet each of you at the door, because each of you is
right, and to each the book gives something different.

If you believe only that a person may believe as they please, you will meet nothing here that tries to bully
you out of it. The instrument compels no belief. What it does, patiently, is separate two things most accounts
leave tangled: the features of a record that are facts about the world from the features that are facts about
the bookkeeping we lay over it --- the marks the world fixes from the marks that are ours to set. Rename,
re-center, re-clock the whole account, and whatever stands unmoved is the physics; whatever shifts was never
more than our description. \emph{The labels are ours; the count is the world's.} That sentence is the spine of
the book's method; you will meet it in the first chapter and again, turned around, at the last. But what you
make of the count once it is in --- that the instrument never presumes to say. It hands you the reading and
steps back.

If you believe in logic, in an argument followed wherever it leads, you will find the book interesting, and
the interest is structural. Its argument is sound, and it is built --- on purpose --- to withhold its own
conclusion. It follows a path all the way around a loop and watches that loop return not quite as it left,
carrying a sign it did not set out with; and exactly where you expect it to close, to announce the result and
take its bow, it stops one step short and turns the final step over to you. An argument that earns every step
but the last, and then declines to take it, is not a gap in the reasoning. It is the reasoning's whole point.
The thing it leaves unproven, it leaves unproven because the proof was never the party entitled to settle it.

If you believe in the laws of science --- believe, fully, that they govern, without favor or exception --- then
you should be warned that to you this book will be cruel, and it will not soften the fact. The laws bind. No
run of sunrises, however unbroken, compels the next: that the sun will rise again is never handed over by the
past, only assumed. Heat will not climb back into the cup it left. The walls come faster as the book goes on
--- no signal outrunning light, no record curling back to before itself, nothing read finer than the floor, no
accounting of the world that lowers its own disorder. To one who holds that the laws govern all, there is no
exemption written anywhere, and no court above them to hear an appeal. Determinism is a cage. This book will
not quietly file the bars down to spare your feelings; it will show them to you straight, and let them hold.

And if you believe in God, there is a place for you here too --- not a proof, but a room. For all the cruelty
of its laws, the argument is built to withhold one thing, deliberately: it will not fix which way you face. It
shuts every door it can reach, and the single door it cannot shut is the one long said to open onto will, or
meaning, or the divine. The book does not claim that anything stands in that room; it carries no instrument
that reaches there. It claims only this --- that the laws, having taken everything else, cannot take that, and
that it will not clear the room out on their behalf. An honest instrument names the floor it stands on and
reports nothing beneath it. This is the floor. Below it the book does not pretend to see, and it will not
pretend the room is empty merely because it cannot see in.

That room --- the one thing the instrument measures everything around and yet cannot itself fix --- has a name.
I will set it down once, here, and then put it away for the length of the book: free will.

What follows is a physics book and nothing else. Past this page it speaks only of counts and clocks, of fields
and forces, of the quantum's blur and the slow bending of light past a star, and it does not say the word again.
You are meant to forget I said it --- to fall into the machinery, to argue with the figures, to lose entirely
the thread of why any of this should bear on you. The word goes under, and stays under. But it does not leave.
It waits in every chapter, beneath the one line the instrument keeps drawing --- between what the world fixes
and what is left to you --- and a reader who reaches the end and turns back to the beginning will find it had
been in plain view the whole way down. Nothing will have been added on the second reading. Something will only
have stopped hiding it.

At the very end, when the instrument has measured everything it is able to measure, it will hand the word back
to you --- not as a result it derived, for it cannot, but as a verdict only you are placed to return. The
oracle will have done its part. It will have told you, as precisely as a finite thing can manage, what is.
What to do with that was never its to say.

You are free. Begin.
